\section{Comparación entre arquitecturas}
\label{sec:comparacion}

En \loso{}, \eegnet{}, el Conformer y LaBraM \cite{jiang2024labram} (variante Base, con ajuste fino
completo; sección~\ref{sec:representaciones}) no muestran diferencias significativas entre sí y son
equivalentes dentro de $\pm 0.03$, mientras que EEGMamba queda por debajo de las tres con tamaño de
efecto grande. En el protocolo intra-sujeto ninguna comparación es equivalente y las diferencias
dependen del procedimiento de selección del modelo. La tabla~\ref{tab:conformer} contrasta las
cuatro arquitecturas, entrenadas con el esquema unificado y evaluadas sobre las mismas $54$
particiones.

\begin{table}[H]
\centering
\caption{\eegnet{}, Conformer, LaBraM (variante Base) y EEGMamba sobre las mismas 54
particiones, pareado por sujeto, con la media de tres semillas por sujeto. El \dsi{} se calcula como
$(m_{\text{visto}} - m_{\text{no visto}})/m_{\text{visto}}$ sobre las medias
agregadas de los 54 sujetos. El símbolo
$\pm$ es la desviación típica entre los 54 sujetos; para los sujetos vistos es la
desviación entre los 54 \textit{folds}, cada uno con unos 51 participantes, de modo
que no mide dispersión entre personas. El kappa de los sujetos vistos es $2\,\mathrm{exactitud} - 1$.
El valor intra-sujeto de EEGMamba emplea la cabeza que conserva la representación por
canal (sección~\ref{sec:eegmamba-adapt}).}
\label{tab:conformer}
\footnotesize
\setlength{\tabcolsep}{3pt}
\begin{tabular}{lcccc}
\toprule
 & \textbf{\eegnet{}} & \textbf{Conformer} & \textbf{LaBraM} & \textbf{EEGMamba} \\
\midrule
Exactitud intra-sujeto & $0.634 \pm 0.107$ & $0.602 \pm 0.117$ & $0.663 \pm 0.153$ & $0.571 \pm 0.095$ \\
\midrule
Exactitud \loso{}, sujetos vistos & $0.781 \pm 0.014$ & $0.775 \pm 0.032$ & $0.796 \pm 0.010$ & $0.742 \pm 0.009$ \\
Exactitud \loso{}, sujeto no visto & $0.761 \pm 0.119$ & $0.747 \pm 0.121$ & $0.755 \pm 0.107$ & $0.715 \pm 0.103$ \\
\dsi{} (exactitud) & $+0.026$ & $+0.037$ & $+0.052$ & $+0.036$ \\
\midrule
Kappa \loso{}, sujetos vistos & $0.561 \pm 0.029$ & $0.550 \pm 0.064$ & $0.593 \pm 0.020$ & $0.484 \pm 0.019$ \\
Kappa \loso{}, sujeto no visto & $0.521 \pm 0.237$ & $0.494 \pm 0.242$ & $0.510 \pm 0.214$ & $0.431 \pm 0.207$ \\
\dsi{} (kappa) & $+0.071$ & $+0.103$ & $+0.140$ & $+0.110$ \\
\bottomrule
\end{tabular}
\end{table}

Bajo \loso{}, \eegnet{}, el Conformer y LaBraM \textbf{no muestran diferencias
significativas} entre sí (tamaños de efecto $r$ de $0.11$ a $0.24$, tabla~\ref{tab:wilcoxon}) y son
equivalentes dentro de $\pm 0.03$ (sección~\ref{sec:wilcoxon}). LaBraM presenta el \dsi{} más alto
(\dsi{} kappa $+0.140$), sin que la diferencia sea concluyente. EEGMamba, la única arquitectura de espacio
de estados, alcanza $0.715$ sobre el sujeto no visto y se separa de las otras tres de forma significativa
y con tamaño de efecto grande, con un \dsi{} del orden del de las demás ($+0.036$).

En el protocolo intra-sujeto, \eegnet{} y LaBraM superan al Conformer de forma significativa
(\eegnet{}: $p < 0.001$, $r = 0.59$; LaBraM: $p = 0.001$, $r = 0.44$), y entre ellas no se detecta
diferencia ($p = 0.44$, $r = 0.11$) pero tampoco puede declararse equivalencia, porque la diferencia de
medias ($+0.03$ a favor de LaBraM, intervalo al $90$\,\% de $[+0.001,\ +0.059]$) sale del margen y la mediana
favorece a \eegnet{}. EEGMamba alcanza $0.571$, la cifra más baja pero por encima del azar
($p = 1.2\times10^{-9}$). Estas diferencias no deben leerse como una comparación de capacidad ni de
preentrenamiento: el criterio de parada conserva un modelo de las épocas de calentamiento en entre el
$53.1$ y el $60.3$\,\% de las ejecuciones del Conformer, frente a entre el $12.4$ y el $16.8$\,\% de
\eegnet{} y entre el $21.1$ y el $28.1$\,\% de LaBraM (sección~\ref{sec:seleccion-intra}).

\subsection{Contrastes pareados y tamaño de efecto}
\label{sec:wilcoxon}

En \loso{}, solo EEGMamba se separa de las demás, con $r$ de $0.54$ a $0.74$; en el protocolo
intra-sujeto se separan el Conformer y EEGMamba. La tabla~\ref{tab:wilcoxon} recoge las seis
comparaciones por protocolo (Wilcoxon pareado, $n = 54$) con el tamaño de efecto
$r = |Z|/\sqrt{n}$ (referencia habitual: $0.1$ pequeño, $0.3$ mediano, $0.5$ grande).

\begin{table}[H]
\centering
\caption{Contrastes de Wilcoxon pareados sobre los 54 participantes, con la media de las tres
semillas de cada participante, y tamaño de efecto $r = |Z|/\sqrt{n}$. La diferencia
mediana corresponde a la primera arquitectura menos la segunda. Cada protocolo y cada métrica
forman una familia de seis contrastes; con la corrección de Holm (exactitud) ninguna conclusión de
significación cambia. Los valores $p$ son bilaterales. En negrita, $r \geq 0.5$. Los resultados por
semilla figuran en la tabla~\ref{tab:variabilidad-dif}.}
\label{tab:wilcoxon}
\small
\begin{tabular}{lrrrrrrr}
\toprule
& \multicolumn{4}{c}{\textbf{Exactitud}} & \multicolumn{3}{c}{\textbf{Kappa}} \\
\cmidrule(lr){2-5}\cmidrule(lr){6-8}
\textbf{Comparación} & \textbf{dif.} & $p$ & $p_{\text{Holm}}$ & $r$ & \textbf{dif.} & $p$ & $r$ \\
\midrule
\multicolumn{8}{l}{\textit{Inter-sujeto (\loso{})}} \\
\eegnet{} vs Conformer & $+0.006$ & 0.073 & 0.22 & $0.24$ & $+0.013$ & 0.073 & $0.24$ \\
\eegnet{} vs LaBraM & $+0.002$ & 0.42 & 0.42 & $0.11$ & $+0.004$ & 0.42 & $0.11$ \\
Conformer vs LaBraM & $-0.008$ & 0.19 & 0.38 & $0.18$ & $-0.017$ & 0.19 & $0.18$ \\
\eegnet{} vs EEGMamba & $+0.045$ & $<0.001$ & $<0.001$ & $\mathbf{0.73}$ & $+0.091$ & $<0.001$ & $\mathbf{0.73}$ \\
Conformer vs EEGMamba & $+0.036$ & $<0.001$ & $<0.001$ & $\mathbf{0.54}$ & $+0.072$ & $<0.001$ & $\mathbf{0.54}$ \\
LaBraM vs EEGMamba & $+0.039$ & $<0.001$ & $<0.001$ & $\mathbf{0.74}$ & $+0.078$ & $<0.001$ & $\mathbf{0.74}$ \\
\midrule
\multicolumn{8}{l}{\textit{Intra-sujeto}} \\
\eegnet{} vs Conformer & $+0.033$ & $<0.001$ & $<0.001$ & $\mathbf{0.59}$ & $+0.067$ & $<0.001$ & $\mathbf{0.59}$ \\
\eegnet{} vs LaBraM & $+0.018$ & 0.44 & 0.44 & $0.11$ & $+0.037$ & 0.44 & $0.11$ \\
Conformer vs LaBraM & $-0.034$ & 0.001 & 0.002 & $0.44$ & $-0.068$ & 0.001 & $0.44$ \\
\eegnet{} vs EEGMamba & $+0.064$ & $<0.001$ & $<0.001$ & $\mathbf{0.73}$ & $+0.127$ & $<0.001$ & $\mathbf{0.73}$ \\
Conformer vs EEGMamba & $+0.027$ & $<0.001$ & $<0.001$ & $\mathbf{0.54}$ & $+0.053$ & $<0.001$ & $\mathbf{0.54}$ \\
LaBraM vs EEGMamba & $+0.044$ & $<0.001$ & $<0.001$ & $\mathbf{0.80}$ & $+0.088$ & $<0.001$ & $\mathbf{0.80}$ \\
\bottomrule
\end{tabular}
\end{table}

\eegnet{}, el Conformer y LaBraM son \textbf{equivalentes dentro del margen de $\pm 0.03$} en
\loso{} (tabla~\ref{tab:equivalencia}): con la media de tres semillas, los intervalos al $90$\,\% de
sus diferencias son $[+0.004,\ +0.025]$ (\eegnet{} y Conformer), $[-0.003,\ +0.015]$ (\eegnet{} y
LaBraM) y $[-0.018,\ +0.002]$ (Conformer y LaBraM), de modo que la equivalencia valdría para cualquier
margen superior a $0.025$. Tres salvedades la acotan: el margen se fijó tras un primer análisis de los
datos y es del orden de la diferencia entre vistos y no vistos ($0.02$ a $0.04$), por lo que no excluye
diferencias de la magnitud del efecto que estudia el \dsi{}; las semillas no comparten régimen de
optimización, aunque la equivalencia se mantiene con las semillas $2027$ y $2028$ solas
(sección~\ref{sec:confusor}); y la dispersión entre semillas de las diferencias, de $0.001$ a $0.006$,
se estima con tres valores (sección~\ref{sec:variabilidad-semillas}).

EEGMamba no es equivalente a ninguna de las otras tres: sus intervalos al $90$\,\% son
$[+0.035,\ +0.056]$ frente a \eegnet{}, $[+0.018,\ +0.044]$ frente al Conformer y
$[+0.030,\ +0.048]$ frente a LaBraM, y la separación se mantiene con las semillas $2027$ y $2028$ solas.

En el protocolo intra-sujeto el Conformer difiere de \eegnet{} y de LaBraM con efectos grande y mediano
($r = 0.59$ y $0.44$) y ningún contraste cumple la equivalencia, con la salvedad de la
sección~\ref{sec:seleccion-intra}. EEGMamba entra en esos contrastes con la cabeza \texttt{flatten},
con la que alcanza $0.571$, frente a una exactitud próxima al azar con la cabeza de promediado global
($0.519$, semilla 2026); queda por debajo de las otras tres y es la más afectada por la selección del
modelo.

\subsubsection{Selección del modelo en el protocolo intra-sujeto}
\label{sec:seleccion-intra}

Con $36$ ensayos de validación, frente a los $800$ de \loso{}, el criterio de parada del protocolo
intra-sujeto conserva con frecuencia un modelo de las diez épocas de calentamiento (parada en la
época $30$, el mínimo del esquema): en el Conformer en más de la mitad de las ejecuciones y en
EEGMamba en casi todas (tabla~\ref{tab:seleccion-intra}).

\begin{table}[H]
\centering
\caption{Proporción de ejecuciones en que el modelo conservado procede de las diez épocas de
calentamiento, por semilla y protocolo. Cada proporción es sobre 54 ejecuciones en \loso{} y
sobre 1\,080 en el protocolo intra-sujeto.}
\label{tab:seleccion-intra}
\small
\begin{tabular}{lcccccc}
\toprule
 & \multicolumn{3}{c}{\textbf{Intra-sujeto}} & \multicolumn{3}{c}{\textbf{\loso{}}} \\
\cmidrule(lr){2-4}\cmidrule(lr){5-7}
 & 2026 & 2027 & 2028 & 2026 & 2027 & 2028 \\
\midrule
\eegnet{} & 16.6\,\% & 16.8\,\% & 12.4\,\% & 7.4\,\% & 7.4\,\% & 3.7\,\% \\
Conformer & 53.1\,\% & 59.3\,\% & 60.3\,\% & 13.0\,\% & 25.9\,\% & 18.5\,\% \\
LaBraM & 21.1\,\% & 28.1\,\% & 24.8\,\% & 25.9\,\% & 33.3\,\% & 24.1\,\% \\
EEGMamba & 85.6\,\% & 88.6\,\% & 86.3\,\% & 46.3\,\% & 46.3\,\% & 40.7\,\% \\
\bottomrule
\end{tabular}
\end{table}

\begin{table}[H]
\centering
\caption{Exactitud media de las ejecuciones cuyo modelo conservado procede de las diez épocas de
calentamiento (Esas, con su número entre paréntesis) y de las demás (Resto), con las tres semillas
agrupadas: $3\,240$ ejecuciones por arquitectura en intra-sujeto y $162$ en \loso{}. La unidad es la
ejecución y no el participante; la comparación es descriptiva.}
\label{tab:seleccion-exactitud}
\small
\begin{tabular}{lcccc}
\toprule
 & \multicolumn{2}{c}{\textbf{Intra-sujeto}} & \multicolumn{2}{c}{\textbf{\loso{}}} \\
\cmidrule(lr){2-3}\cmidrule(lr){4-5}
 & Esas & Resto & Esas & Resto \\
\midrule
\eegnet{}  & 0.511 (494)  & 0.656 & 0.724 (10) & 0.763 \\
Conformer  & 0.521 (1865) & 0.714 & 0.685 (31) & 0.761 \\
LaBraM     & 0.521 (800)  & 0.710 & 0.707 (45) & 0.773 \\
EEGMamba   & 0.541 (2814) & 0.767 & 0.703 (72) & 0.725 \\
\bottomrule
\end{tabular}
\end{table}

En esas ejecuciones la exactitud media intra-sujeto es de $0.51$ a $0.54$, próxima al azar, frente a
$0.66$ a $0.77$ en las restantes (tabla~\ref{tab:seleccion-exactitud}); la comparación es descriptiva,
porque las ejecuciones que prolongan el entrenamiento pueden corresponder a participantes más
favorables. Por tanto, la exactitud intra-sujeto del Conformer y de EEGMamba refleja en buena parte el
comportamiento del criterio de selección con $36$ ensayos y no solo la capacidad de aprender de $144$,
y las diferencias intra-sujeto entre arquitecturas no pueden atribuirse a la capacidad ni al
preentrenamiento. En \loso{} el fenómeno es menos frecuente (la mayor proporción es la de EEGMamba,
de $40.7$ a $46.3$\,\%) y la exactitud de las ejecuciones afectadas difiere poco de la del resto
($0.703$ frente a $0.725$).

\subsubsection{Número de semillas por arquitectura}
\label{sec:semillas-por-arq}

Las cuatro arquitecturas se ejecutaron con tres semillas ($2026$, $2027$ y $2028$) en los protocolos
intra-sujeto, \loso{} y de transferencia, lo que da $54 \times 20 \times 3 = 3\,240$ mediciones
intra-sujeto y $54 \times 3 = 162$ \textit{folds} \loso{} por arquitectura
(sección~\ref{sec:semillas-metodo}). El régimen de optimización no es el mismo en todas: \eegnet{}, el
Conformer y LaBraM tienen la semilla $2026$ sin recorte de gradiente y las otras dos con él, salvo en el
protocolo intra-sujeto de \eegnet{} y del Conformer, donde las tres se entrenaron sin recorte
(tabla~\ref{tab:regimen}); la sección~\ref{sec:confusor} acota cuánto puede haber contribuido ese cambio.

La mayor dispersión intra-sujeto de LaBraM ($0.153$ frente a $0.107$ de \eegnet{}) procede de la
variabilidad entre participantes y no de las semillas: la desviación entre semillas de la exactitud media
es de $0.001$ a $0.006$, y la de un mismo participante de $0.016$ a $0.025$, frente a $0.095$ a $0.153$
entre participantes. Los contrastes pareados emplean la media de las tres semillas de cada participante
(sección~\ref{sec:estadistica}).

\subsection{Intervalos de confianza por remuestreo}
\label{sec:bootstrap-resultados}

Los intervalos por remuestreo (sección~\ref{sec:bootstrap}, tabla~\ref{tab:bootstrap}) confirman las
conclusiones anteriores: el \dsi{} excluye el cero solo en LaBraM, y la separación de EEGMamba es el
único ordenamiento entre arquitecturas que el remuestreo sostiene.

\begin{table}[htbp]
\centering
\caption{Intervalos de confianza al $95$\,\% por remuestreo jerárquico con reemplazo (participantes
y, dentro de cada participante, sus tres semillas), $B = 10\,000$, semilla $2026$, sobre los $54$
valores por participante, para el \dsi{} y las diferencias entre arquitecturas con la media de tres
semillas. En las diferencias, la extracción de participantes es común a las dos arquitecturas del
par. El intervalo es de percentiles (2.5 y 97.5) de las $B$ réplicas. Las filas de exactitud
(bloque inferior) son la media de tres semillas, con el mismo remuestreo jerárquico.}
\label{tab:bootstrap}
\small
\begin{tabular}{lccc}
\toprule
 & \textbf{Estimación} & \textbf{IC $95$\,\%} & \textbf{Amplitud} \\
\midrule
\multicolumn{4}{l}{\textit{\dsi{}, media de tres semillas}} \\
\eegnet{} & $+0.026$ & $[-0.016,\ +0.068]$ & 0.084 \\
Conformer & $+0.037$ & $[-0.005,\ +0.078]$ & 0.084 \\
LaBraM & $+0.052$ & $[+0.016,\ +0.088]$ & 0.072 \\
EEGMamba & $+0.036$ & $[-0.002,\ +0.074]$ & 0.076 \\
\midrule
\multicolumn{4}{l}{\textit{Diferencia de exactitud \loso{}, media pareada de tres semillas}} \\
\eegnet{} $-$ Conformer & $+0.014$ & $[+0.000,\ +0.029]$ & 0.029 \\
\eegnet{} $-$ LaBraM & $+0.006$ & $[-0.006,\ +0.018]$ & 0.023 \\
Conformer $-$ LaBraM & $-0.008$ & $[-0.022,\ +0.006]$ & 0.028 \\
\eegnet{} $-$ EEGMamba & $+0.045$ & $[+0.032,\ +0.058]$ & 0.026 \\
Conformer $-$ EEGMamba & $+0.032$ & $[+0.014,\ +0.048]$ & 0.034 \\
LaBraM $-$ EEGMamba & $+0.040$ & $[+0.028,\ +0.050]$ & 0.023 \\
\midrule
\multicolumn{4}{l}{\textit{Diferencia de exactitud intra-sujeto, media pareada de tres semillas}} \\
\eegnet{} $-$ Conformer & $+0.032$ & $[+0.018,\ +0.045]$ & 0.027 \\
\eegnet{} $-$ LaBraM & $-0.029$ & $[-0.065,\ +0.005]$ & 0.070 \\
Conformer $-$ LaBraM & $-0.061$ & $[-0.093,\ -0.031]$ & 0.062 \\
\eegnet{} $-$ EEGMamba & $+0.063$ & $[+0.046,\ +0.082]$ & 0.036 \\
Conformer $-$ EEGMamba & $+0.032$ & $[+0.018,\ +0.047]$ & 0.029 \\
LaBraM $-$ EEGMamba & $+0.093$ & $[+0.063,\ +0.123]$ & 0.060 \\
\bottomrule
\end{tabular}

\medskip
\small
\begin{tabular}{lccc}
\toprule
 & \textbf{Estimación} & \textbf{IC $95$\,\%} & \textbf{Amplitud} \\
\midrule
\multicolumn{4}{l}{\textit{Media de tres semillas, remuestreo jerárquico}} \\
\multicolumn{4}{l}{\textit{Exactitud \loso{}}} \\
\eegnet{}  & 0.761 & $[0.728,\ 0.792]$ & 0.064 \\
Conformer  & 0.747 & $[0.714,\ 0.780]$ & 0.066 \\
LaBraM     & 0.755 & $[0.726,\ 0.783]$ & 0.057 \\
EEGMamba   & 0.715 & $[0.688,\ 0.743]$ & 0.055 \\
\midrule
\multicolumn{4}{l}{\textit{Exactitud intra-sujeto}} \\
\eegnet{}  & 0.634 & $[0.606,\ 0.664]$ & 0.058 \\
Conformer  & 0.602 & $[0.572,\ 0.635]$ & 0.062 \\
LaBraM     & 0.663 & $[0.624,\ 0.706]$ & 0.082 \\
EEGMamba   & 0.571 & $[0.547,\ 0.598]$ & 0.050 \\
\bottomrule
\end{tabular}
\end{table}

Las diferencias entre \eegnet{}, el Conformer y LaBraM en \loso{} tienen intervalos de amplitud
$0.023$ a $0.029$ que contienen el cero, salvo el de \eegnet{} frente al Conformer, que lo excluye por
$0.0005$ sin corregir por las seis comparaciones y con $p = 0.073$ en el contraste pareado; no se
declara diferencia. Los intervalos de las tres diferencias que involucran a EEGMamba
($[+0.032,\ +0.058]$, $[+0.014,\ +0.048]$ y $[+0.028,\ +0.050]$) excluyen el cero con holgura. La tabla
reporta diferencias \emph{medias}, y la tabla~\ref{tab:wilcoxon} diferencias \emph{medianas}; pueden
discrepar (entre \eegnet{} y el Conformer, $+0.006$ y $+0.014$) por la asimetría de la distribución por
participante. Al parear por participante la variabilidad entre personas se cancela en buena parte y la
inicialización pasa a ser del $17$ al $35$\,\% de la varianza de la diferencia en \loso{}
(sección~\ref{sec:variabilidad-semillas}).

\subsection{Variabilidad por semilla de las diferencias pareadas}
\label{sec:variabilidad-semillas}

La diferencia media entre arquitecturas tiene el mismo signo en las tres semillas en los $12$
contrastes de exactitud, y ninguna conclusión de significación ni de equivalencia cambia con tres
semillas (tabla~\ref{tab:variabilidad-dif}). Lo que varía entre semillas es el valor $p$ de una semilla
sola y, en los contrastes pequeños de \loso{}, la propia diferencia.

\begin{table}[H]
\centering
\caption{Diferencia media de exactitud (primera menos segunda arquitectura) en cada semilla, desviación
típica (DT) entre las tres semillas de esa diferencia, valores $p$ de Wilcoxon de cada semilla por separado
(2026 / 2027 / 2028) e intervalo al $95$\,\% jerárquico de la diferencia de la media de tres
semillas.}
\label{tab:variabilidad-dif}
\footnotesize
\setlength{\tabcolsep}{3pt}
\resizebox{\textwidth}{!}{%
\begin{tabular}{lccccll}
\toprule
 & \multicolumn{3}{c}{\textbf{Dif. media por semilla}} & & & \\
\cmidrule(lr){2-4}
\textbf{Contraste} & 2026 & 2027 & 2028 & DT & $p$ por semilla & IC $95$\,\% (media de 3) \\
\midrule
\multicolumn{7}{l}{\textit{Inter-sujeto (\loso{})}} \\
\eegnet{} $-$ Conformer & $+0.011$ & $+0.015$ & $+0.015$ & 0.0025 & 0.12 / 0.063 / 0.49 & $[+0.000,\ +0.029]$ \\
\eegnet{} $-$ LaBraM & $+0.009$ & $+0.006$ & $+0.002$ & 0.0035 & 0.27 / 0.41 / 0.88 & $[-0.006,\ +0.018]$ \\
Conformer $-$ LaBraM & $-0.002$ & $-0.009$ & $-0.013$ & 0.0055 & 0.49 / 0.32 / 0.50 & $[-0.022,\ +0.006]$ \\
\eegnet{} $-$ EEGMamba & $+0.046$ & $+0.046$ & $+0.044$ & 0.0008 & $<0.001$ / $<0.001$ / $<0.001$ & $[+0.032,\ +0.058]$ \\
Conformer $-$ EEGMamba & $+0.035$ & $+0.030$ & $+0.029$ & 0.0030 & $<0.001$ / $<0.001$ / $<0.001$ & $[+0.014,\ +0.048]$ \\
LaBraM $-$ EEGMamba & $+0.037$ & $+0.039$ & $+0.042$ & 0.0027 & $<0.001$ / $<0.001$ / $<0.001$ & $[+0.028,\ +0.050]$ \\
\midrule
\multicolumn{7}{l}{\textit{Intra-sujeto}} \\
\eegnet{} $-$ Conformer & $+0.026$ & $+0.031$ & $+0.038$ & 0.0062 & $<0.001$ / $<0.001$ / $<0.001$ & $[+0.018,\ +0.045]$ \\
\eegnet{} $-$ LaBraM & $-0.032$ & $-0.033$ & $-0.023$ & 0.0052 & 0.30 / 0.31 / 0.58 & $[-0.065,\ +0.005]$ \\
Conformer $-$ LaBraM & $-0.058$ & $-0.063$ & $-0.061$ & 0.0028 & 0.004 / 0.001 / $<0.001$ & $[-0.093,\ -0.031]$ \\
\eegnet{} $-$ EEGMamba & $+0.056$ & $+0.060$ & $+0.074$ & 0.0097 & $<0.001$ / $<0.001$ / $<0.001$ & $[+0.046,\ +0.082]$ \\
Conformer $-$ EEGMamba & $+0.030$ & $+0.029$ & $+0.036$ & 0.0038 & $<0.001$ / $<0.001$ / $<0.001$ & $[+0.018,\ +0.047]$ \\
LaBraM $-$ EEGMamba & $+0.088$ & $+0.092$ & $+0.098$ & 0.0049 & $<0.001$ / $<0.001$ / $<0.001$ & $[+0.063,\ +0.123]$ \\
\bottomrule
\end{tabular}
}
\end{table}

La desviación entre semillas de la diferencia media va de $0.0008$ a $0.0055$ en \loso{} y de $0.003$ a
$0.010$ en el protocolo intra-sujeto. En los contrastes pequeños de \loso{} es del orden del propio
efecto (\eegnet{} frente a LaBraM vale $+0.009$, $+0.006$ y $+0.002$ en las tres semillas) y en los de
EEGMamba es despreciable. El valor $p$ de una semilla sola puede cambiar mucho: \eegnet{} frente al
Conformer en \loso{} da $0.12$, $0.063$ y $0.49$, y $0.073$ con la media de las tres. El ruido de semilla
explica del $17$ al $35$\,\% de la varianza entre participantes de la diferencia media en \loso{} y del
$2$ al $23$\,\% en el protocolo intra-sujeto.

No se detecta que el cambio de régimen modifique estas diferencias. El desplazamiento entre la diferencia
de la semilla $2026$ y la de las semillas $2027$ y $2028$ va de $-0.011$ a $+0.009$ en los doce contrastes,
con intervalos al $95$\,\% que contienen el cero, y en el contraste sin cambio de régimen (\eegnet{}
frente al Conformer en intra-sujeto) vale $-0.008$ $[-0.018,\ +0.001]$. La resolución de esta comparación
es de $\pm 0.01$ a $\pm 0.02$, de modo que los datos no distinguen el efecto del régimen del ruido
entre semillas.

En la transferencia, las semillas cambian la estimación puntual y no la inferencia: la desviación entre
semillas de la diferencia entre la caída por protocolo y por persona es de $0.003$ a $0.029$, frente a
una semianchura del intervalo de $0.12$ a $0.16$. Con tres semillas la desviación entre semillas tiene
dos grados de libertad, el emparejamiento de semillas entre arquitecturas solo las etiqueta por índice y
las semillas no son réplicas del mismo entorno (sección~\ref{sec:semillas-por-arq}); la variabilidad por
semilla queda estimada solo de forma aproximada.

\subsection{Equivalencia práctica en ambos protocolos}
\label{sec:equivalencia-resultados}

En \loso{}, \eegnet{}, el Conformer y LaBraM son equivalentes dentro de $\pm 0.03$ y EEGMamba no es
equivalente a ninguna; en el protocolo intra-sujeto ningún par lo es (tabla~\ref{tab:equivalencia}). Como
el margen se fijó tras un primer análisis de los datos, se informa también el menor margen con el que se
cumpliría la equivalencia.

\begin{table}[H]
\centering
\caption{Equivalencia práctica con margen $\pm 0.03$ en exactitud. Entre paréntesis, si el
intervalo al $90$\,\% queda dentro del margen. Margen mín.: mayor valor absoluto de los extremos
del intervalo con la media de tres semillas. $p_{\text{TOST,Holm}}$: contraste TOST con la prueba
de Wilcoxon, corregido por Holm en la familia de seis.}
\label{tab:equivalencia}
\small
\resizebox{\textwidth}{!}{%
\begin{tabular}{lllcc}
\toprule
\textbf{Contraste} & IC $90$\,\% (2026) & IC $90$\,\% (media de 3) & Margen mín. & $p_{\text{TOST,Holm}}$ \\
\midrule
\multicolumn{5}{l}{\textit{Inter-sujeto (\loso{})}} \\
\eegnet{} $-$ Conformer & $[+0.002,\ +0.020]$ (sí) & $[+0.004,\ +0.025]$ (sí) & 0.025 & 0.001 \\
\eegnet{} $-$ LaBraM & $[-0.002,\ +0.020]$ (sí) & $[-0.003,\ +0.015]$ (sí) & 0.015 & $<0.001$ \\
Conformer $-$ LaBraM & $[-0.013,\ +0.009]$ (sí) & $[-0.018,\ +0.002]$ (sí) & 0.018 & 0.003 \\
\eegnet{} $-$ EEGMamba & $[+0.035,\ +0.057]$ (no) & $[+0.035,\ +0.056]$ (no) & 0.056 & 1.00 \\
Conformer $-$ EEGMamba & $[+0.023,\ +0.046]$ (no) & $[+0.018,\ +0.044]$ (no) & 0.044 & 1.00 \\
LaBraM $-$ EEGMamba & $[+0.028,\ +0.046]$ (no) & $[+0.030,\ +0.048]$ (no) & 0.048 & 1.00 \\
\midrule
\multicolumn{5}{l}{\textit{Intra-sujeto}} \\
\eegnet{} $-$ Conformer & $[+0.014,\ +0.038]$ (no) & $[+0.021,\ +0.042]$ (no) & 0.042 & 1.00 \\
\eegnet{} $-$ LaBraM & $[-0.062,\ -0.002]$ (no) & $[-0.059,\ -0.001]$ (no) & 0.059 & 1.00 \\
Conformer $-$ LaBraM & $[-0.085,\ -0.032]$ (no) & $[-0.087,\ -0.036]$ (no) & 0.087 & 1.00 \\
\eegnet{} $-$ EEGMamba & $[+0.040,\ +0.073]$ (no) & $[+0.049,\ +0.078]$ (no) & 0.078 & 1.00 \\
Conformer $-$ EEGMamba & $[+0.017,\ +0.044]$ (no) & $[+0.020,\ +0.044]$ (no) & 0.044 & 1.00 \\
LaBraM $-$ EEGMamba & $[+0.063,\ +0.113]$ (no) & $[+0.068,\ +0.118]$ (no) & 0.118 & 1.00 \\
\bottomrule
\end{tabular}
}
\end{table}

La equivalencia en \loso{} se mantiene con la media de tres semillas, con el intervalo jerárquico y con las
semillas $2027$ y $2028$ solas (todas con recorte de gradiente). El menor margen es de $0.025$ para
\eegnet{} frente al Conformer con tres semillas y de $0.029$ con las semillas $2027$ y $2028$, de modo que
en ese par la conclusión queda justo dentro del margen y no vale para márgenes menores. El TOST con la
prueba de Wilcoxon la confirma ($p_{\text{Holm}}$ de $0.001$, $<0.001$ y $0.003$), y el TOST con la $t$ de
Student, más conservador, también. Con EEGMamba los márgenes mínimos son de $0.044$ a $0.056$.

En el protocolo intra-sujeto \textbf{ningún contraste es equivalente}: el menor margen va de $0.042$ a
$0.118$ con tres semillas. \eegnet{} frente a LaBraM, sin diferencia detectable ($p = 0.44$), tampoco es
equivalente, de modo que no puede declararse ni diferencia ni igualdad.

